\scanpage{59}
Ainsi $\unsure{\psi|\varphi}$ \uline{est un $1$-cocycle}, quand on fait opérer $\operatorname{Gal}(k/k_0)$ sur $G(k)$ de la façon évidente (par conjugaison). Il correspond à un fibré principal homogène $P_{v,k}$ sous $G(k)$, avec opérations de $\operatorname{Gal}(k/k_0)$ sur \deleted{$(P_v,G(k))$ compatibles avec le grp}. Il définit \inserted{par descente} un fibré principal homogène \inserted{$P_{v,k_0}$} sous $G_{k_0}$, avec un \uline{splitting donné} sur $k$.

On peut d'ailleurs dire ce que c'est que $P$ : c'est le fibré principal des isom \inserted{de foncteurs} $\xi_{M,\unsure{k_0|\infty}}$ avec $\xi_{M,v}^{\mathbf Z}\otimes_{\mathbf Z}k_0$, dont une trivialisation est en effet obtenue $\ldots$
\margin{\uline{À vérifier}\\OK.}

\unsure{Pour|Puis} la conjecture \inserted{de transcendance} N° 1, \deleted{\unsure{celle que l'hom}} \unsure{ce|le} point ainsi donné sur $P_{v,k}$ est \uline{générique}, revient à dire que l'application
\diagram{\begin{tikzcd}[ampersand replacement=\&,column sep=large]
\operatorname{Gal}(k/k_0)\arrow[r,"\varphi_v"]\&G(k)[\subset\mathrm{Gl}(M_k(\infty))]
\end{tikzcd}}
\deleted{est} a une image Zariski dense. [NB. elle l'est \unsure{déjà},
\margin{$\pi$ transcendant sur $k_0\ldots$}
\mspage{58}
